\section{DISCUSSION}

\ourmethod{} incurs additional prompting overhead due to backward slicing and constraint hint construction. On average, each focal method consumes approximately \textbf{43,000 tokens}, reflecting both the input code context and the appended constraints. Across all 3,576 requests in our evaluation, this corresponds to roughly \textbf{1.43 billion input tokens} and \textbf{12.5 million output tokens}, for a total of about \textbf{1.55 billion tokens}. The backward slicing step helps mitigate this cost by removing irrelevant statements, reducing the prompt length by roughly \textbf{25\%} before constraint hints are appended. In absolute terms, a project with 130 focal methods costs about \textbf{\$100}, which is reasonable considering the resulting coverage improvements.

\textbf{Token Efficiency.} Despite the large aggregate token consumption, \ourmethod{} achieves a favorable coverage-per-token ratio by focusing the LLM on relevant code slices and explicit constraints. This ensures that each token contributes meaningfully to exercising complex branches, improving the overall efficiency of test generation compared to naive full-context prompts.

\textbf{Qualitative Insight.} To illustrate why \ourmethod{} outperforms baseline approaches, we examine a representative example from the \textit{Favicon for Alibaba} project. The method \texttt{generateFavicon} contains a target branch condition \texttt{if (icon != null \&\& icon.size() > 0)}, which can only be satisfied if the object is properly initialized beforehand. Baselines generate tests that invoke the focal method directly, failing to establish the required state and leaving the branch uncovered. \ourmethod{}, by applying backward slicing, identifies the relevant state-setting call \texttt{setIconState} and provides a constraint hint specifying the precondition. The resulting test correctly initializes the object state and triggers the target branch. This example demonstrates that \ourmethod{}’s benefits extend beyond quantitative gains: it enhances the LLM's capacity to reason about cross-method dependencies in realistic Java programs.

\textbf{Limitations.} Despite its effectiveness, \ourmethod{} has notable constraints. \emph{Slice granularity.} For projects such as \textit{Collections} and \textit{Compress}, where branches depend on simple value-range conditions rather than cross-method state, the backward slice may remove auxiliary type or range information, reducing test quality. Conversely, overly coarse slices may retain irrelevant statements, introducing noise. Developing adaptive slicing strategies that adjust granularity based on branch type remains future work. \emph{Scalability.} Backward slicing and constraint extraction introduce preprocessing overhead proportional to the inter-procedural dependency graph. In our evaluation, average preprocessing time per focal method is \textbf{1.2\,s}, acceptable for medium-scale projects, though on-demand or incremental slicing may be required for industrial-scale codebases.